\documentclass{article}
\usepackage[letterpaper,top=2cm,bottom=2cm,left=3cm,right=3cm,marginparwidth=1.75cm]{geometry}

\usepackage{graphicx} % Required for inserting images
\usepackage{booktabs}
\usepackage{caption}
\usepackage{subcaption}
\usepackage{authblk}
\usepackage{amsmath}
\usepackage{amssymb}
\usepackage[dvipsnames]{xcolor}
\usepackage[english]{babel}
\usepackage{pgfgantt}
\usepackage{csquotes}
\usepackage{multirow}

\usepackage[colorlinks=true, allcolors=blue]{hyperref}
\hypersetup{colorlinks=true, breaklinks=true}

\usepackage[    
    style=numeric-comp,
    sorting=none,
    sortcites=true
    ]{biblatex}

\bibliography{latex/bibliography}

\newcommand{\mycitethis}{\tkDE{[?]}}

\usepackage{lscape}

\usepackage{verbatim}
\usepackage{geometry}


% Change font
\usepackage{helvet}
\renewcommand{\familydefault}{\sfdefault}

\usepackage{xcolor}
\newcommand{\tkDE}[1]{\textcolor{magenta}{#1}} 
\newcommand{\citethis}{\tkDE{[?]}}

\newcommand{\diff}{\mathrm{d}}  

\title{BSol Measles Participatory System Modelling - Transcript Extraction}
\author{David Ellis}
\date{}

\begin{document}

\maketitle 

The following is a summary of the recorded transcript of a workshop held on the 20th of March 2026. This workshop, hosted by Birmingham City Council, brought together local experts to build a conceptual model of the local measles system, to identify behavioural and epidemiological factors that facilitate and block flows in the system, and to identify potential interventions to improve these flows.
Due to the length of the workshop, Copilot was used to extract vocal contributions to supplement the written materials. The full Copilot output, can be viewed on  \href{https://github.com/David-Ellis/Measles-System-Dynamics-BSol/blob/main/data/workshop/chronological_extraction_FULL.txt}{GitHub}\footnote{All non-sensitive materials related to this project can also be accessed via the GitHub repository \url{https://github.com/David-Ellis/Measles-System-Dynamics-BSol}.}. 

\section{Copilot Prompt}

The full prompt given to Copilot was as follows:

\begin{verbatim}
I will upload two MS Teams transcript files (Part 1 and Part 2).

I need you to produce a full chronological extraction of contributions from both transcripts, 
ordered strictly by agenda items. The extractions should not include any names or otherwise 
identifiable information. Use the following agenda headings, in this order:

Introduction
10:00 Arrival
10:10 Welcome and project overview
10:20 Local epidemiology update (UKHSA)
10:30 Summary of key evidence on vaccination barriers
11:00 Brief introduction to system dynamics modelling
11:10 System dynamics interactive example (group task)
11:30 Break

Model Building
11:45 Base model introduction
11:55 Identify factors that influence key flows in the system (group task)
12:20 Discussion
12:30 Vote on most influential factors
12:45 Lunch Break

Interventions
13:30 Identify problematic flows in the system
13:40 Identify potential interventions to improve system behaviour (group task)
14:10 Discussion
14:30 Rank proposed interventions using an anonymous online vote
15:00 Finish

Instructions:

1. Go through both transcript files and extract *all* contributions in *chronological order*, 
rewriting them cleanly in full sentences.
2. Place each contribution under the correct agenda heading.
3. If the transcript text is unclear, truncated, or garbled, insert "[UNCLEAR]" at that 
point and summarise what is recoverable.
4. Do not omit any substantive content—include all contributions, group discussions, side 
comments, and facilitator guidance.
5. Produce the final output as a set of .txt files:
   - Part 0 (header/explanation)
   - Part 1 (Introduction section)
   - Part 2 (Model Building section)
   - Part 3 (Interventions section)
6. After generating Parts 1–3, compile them into a single combined file called 
"chronological_extraction_FULL.txt".

Please ask me if you need confirmation about formatting, but otherwise proceed without 
repeatedly checking after each step. Proceed without citations for transcript processing.
\end{verbatim}

\section{Chronological Transcript Extraction}

\subsection{Introduction}
\subsubsection{Welcome}
\begin{itemize}
    \item The facilitator outlined the purpose of the workshop as supporting system dynamics modelling related to measles vaccination uptake and transmission.
    \item Participants were informed that the session would combine evidence review, interactive modelling exercises, and group discussion.
    \item The importance of capturing complexity without over‑complicating the model was emphasised.
    \item The session was framed as supporting both academic work and practical system decision‑making.
    \item The overall project aim was described as developing a system dynamics model of measles transmission and vaccination uptake.
    \item The model was described as a tool to explore how different factors interact rather than a precise forecasting model.
    \item Participants were informed that the focus would be on understanding drivers, feedback loops, and leverage points in the system.
    \item The agenda for the day was briefly outlined.
    \item The importance of participant expertise and open discussion was highlighted.
\end{itemize}

\subsubsection{Local epidemiology update} 

\begin{itemize}
    \item An overview of recent measles vaccination uptake was presented, focusing on first‑dose MMR coverage.
    \item It was noted that first‑dose uptake data are more robust than second‑dose data.
    \item Participants discussed that having a long gap between first and second doses reduces effectiveness.
    \item It was stated that, for modelling simplicity, the initial model would focus on single‑dose uptake.
    \item Recent trends showed a temporary increase in vaccination activity followed by stabilisation.
    \item Birmingham was described as having a younger population than the national average, increasing vaccination demand.
    \item Population diversity and high birth rates were highlighted as important contextual factors.
    \item Participants noted that Birmingham requires proportionally more vaccinations per practice than many other areas.
\end{itemize}

\subsubsection{Summary of key evidence on vaccination barriers}

\begin{itemize}
    \item A summary of barriers to MMR vaccination was presented based on local qualitative intelligence.
    \item Language barriers were identified as a major issue, including limited English literacy and low literacy in any language.
    \item Difficulties accessing and understanding written communications were highlighted.
    \item Concerns about vaccine ingredients, particularly pork gelatine, were discussed.
    \item It was noted that awareness of pork‑free vaccine options varies across communities.
    \item Fears linking MMR vaccination to autism were described as persistent despite being discredited.
    \item Differences in concerns between ethnic and community groups were highlighted.
    \item Migrant and asylum‑seeking populations were described as facing barriers related to access, eligibility awareness, and missing records.
    \item Fear of cost and immigration enforcement was identified as an unnecessary but persistent barrier.
    \item Gypsy, Roma and Traveller communities were described as facing unique access challenges and historical mistrust.
    \item General mistrust of government and institutions was discussed as a background factor.
    \item Participants noted that reasons for vaccine hesitancy are highly individual and rarely explained by a single factor.
    \item It was suggested that reducing the need for individual decision‑making could increase uptake.
    \item (UNCLEAR) Partial discussion referencing national‑level trust and inequality.
\end{itemize}

\subsubsection{Brief introduction to system dynamics modelling}

\begin{itemize}
    \item System dynamics was introduced as a method for modelling complex systems over time.
    \item Key concepts of stocks, flows, and feedback loops were explained.
    \item Stocks were described as accumulations, such as populations or quantities.
    \item Flows were described as rates of change between stocks.
    \item Information flows and abstract variables such as trust were described as model‑able.
    \item The concept of system boundaries and exclusions was explained.
    \item Participants were reminded that models simplify reality and rely on assumptions.
    \item The idea of tuning model parameters to reproduce observed patterns was introduced
\end{itemize}

\subsubsection{System dynamics interactive example (group task)}

\begin{itemize}
    \item A simple example using population dynamics was presented to illustrate stocks and flows.
    \item The example demonstrated how feedback loops can create unintended behaviour.
    \item Participants discussed how missing feedback can lead to unrealistic model behaviour.
    \item The importance of bidirectional relationships was emphasised.
    \item Visual cues such as reinforcing and balancing loops were explained.
    \item Participants were reassured that perfection was not required for the exercise.
    \item Groups were instructed to work together on a simple illustrative task.
    \item (UNCLEAR) Light‑hearted comments and informal discussion during the exercise.
    \item Participants were invited to take a short break.
\end{itemize}


\subsection{Model Building}

\subsubsection{Base model introduction}

\begin{itemize}
    \item The facilitator presented the simplified base model for measles dynamics consisting of five key stocks: susceptible, vaccinated, infected, hospitalised, and recovered.
    \item It was explained that deaths from measles would be included, but other deaths or age transitions would not be modelled for simplicity.
    \item The facilitator noted that although ethnicity was not shown visually, each stock would be broken down by broad ethnic categories to reflect differential uptake and risk.
    \item Participants were invited to add additional stocks if they believed essential system components were missing.
    \item The facilitator highlighted that vaccinated individuals could still, in rare cases, become infected and instructed participants to add that relationship to their diagrams.
    \item Migration‑related factors were discussed, with participants noting that some individuals may enter the system already vaccinated or unvaccinated.
    \item The facilitator encouraged inclusion of new inflows for vaccinated and unvaccinated migrants as needed.
    \item Participants raised questions about recovery categories, including long‑term complications, and how these might be represented.
    \item The facilitator suggested that complication pathways could be included by splitting outflows or by adding abstract indicators linked to recovery.
    \item Participants were reminded that the model should balance simplicity with capturing key system behaviour.
\end{itemize}

\subsubsection{Identify factors that influence key flows in the system (group task)}

\begin{itemize}
    \item Participants were instructed to identify all factors influencing vaccination rate, infection rate, recovery rate, and hospitalisation rate.
    \item Groups were encouraged to write each factor on sticky notes and draw arrows showing causal relations.
    \item The facilitator emphasised capturing abstract drivers, such as trust, perceived risk, and misinformation.
    \item Participants discussed parental willingness to vaccinate as a central factor influencing vaccination rate.
    \item Accessibility of vaccination services was identified as equally critical, including appointment availability, opening hours, and location.
    \item  Participants noted that willingness and accessibility act as parallel necessities: both must be present for vaccination to occur.
    \item Groups identified GP motivation and practice‑level prioritisation as influential factors.
    \item Workforce capacity was highlighted, including staffing levels, availability of vaccinators, and competing pressures.
    \item Participants discussed cost‑effectiveness considerations within practices, influencing how actively staff promote vaccination.
    \item Family influence was identified as a major factor, particularly when relatives have experienced negative side effects or misinformation.
    \item Perceived disease risk was highlighted, with participants noting that fear of outbreaks can both increase and decrease uptake depending on how risk is interpreted.
    \item Social and community norms were identified as important determinants of willingness.
    \item Socioeconomic conditions and deprivation were discussed as constraints on time, transport, and access.
    \item National‑level policy and messaging were noted as background influences outside local control.
    \item Participants discussed the influence of media exposure, including positive messaging, negative messaging, and misinformation spread.
    \item Health literacy was identified as a key determinant affecting understanding of eligibility, safety, and access pathways.
    \item Participants discussed that some communities experience mistrust of institutions, which reduces willingness to engage.
    \item The number of infectious individuals in the community was linked to perceived risk and thus willingness.
    \item The presence of hospitalised cases was noted as potentially influencing population risk perception.
    \item Participants reviewed stock‑and‑flow relationships to ensure each factor pointed to a measurable rate or mechanism.
    \item (UNCLEAR) Partial discussion referencing perceived community risk and outbreak triggers.
\end{itemize}

\subsubsection{Flow Factors Group Discussion}

\begin{itemize}
    \item The whole group reconvened to share the factors identified during the group task.
    \item Participants agreed that accessibility emerged as one of the most influential drivers across groups.
    \item Willingness or intent to vaccinate was identified as the other dominant driver.
    \item It was noted that almost all secondary factors could be linked to one of these two overarching domains.
    \item Participants discussed the relationship between GP incentives and vaccination activity, including how payment structures influence motivation.
    \item It was suggested that practices sometimes deprioritise vaccination if targets appear unattainable.
    \item Participants emphasised that socioeconomic pressures strongly influence both access and willingness.
    \item Several participants highlighted that national messages, policy decisions, or negative media activity influence local behaviour.
    \item Participants raised the issue of cultural and linguistic barriers affecting understanding of vaccine options, especially pork‑free alternatives.
    \item It was noted that fear of side effects, both real and perceived, continues to shape parental decisions.
    \item Participants discussed the significant role of family members and peer networks in shaping beliefs about vaccines.
    \item Some participants observed that fear of disease spread can prompt increased vaccination, while excessive fear may cause avoidance.
    \item Participants noted that schools, community centres, and trusted local figures can influence uptake through advocacy.
    \item It was recognised that time constraints, transport limitations, and clinic scheduling create structural barriers to accessing vaccination.
    \item Participants identified misinformation as a key influence, including the volume and perceived credibility of misleading content.
    \item Data limitations were discussed, including limited visibility of ethnicity‑specific patterns for childhood vaccinations.
    \item Participants noted that uptake disparities reflect both structural inequities and differing community experiences with healthcare.
    \item Several participants discussed that health literacy influences correct interpretation of side effects versus disease symptoms.
    \item Participants affirmed that the model should remain usable rather than overly complex, even if many factors appear relevant.
    \item (UNCLEAR) Fragmented comments about grouping factors and removing duplicates.
\end{itemize}

\subsubsection{Vote on most influential factors}

\begin{itemize}
    \item Participants were given stickers to vote on the three factors they believed most strongly drive system behaviour.
    \item The facilitator encouraged people to review all groups’ notes before voting.
    \item Participants were reminded that multiple votes could be allocated to a single factor if desired.
    \item Accessibility received substantial interest, reflecting strong group agreement on its importance.
    \item Willingness or intention to vaccinate also attracted many votes, confirming its central role.
    \item GP motivation and prioritisation received discussion, though fewer votes than accessibility or willingness.
    \item Health literacy received mixed but notable attention as a moderating factor.
    \item Misinformation and negative messaging were highlighted as widespread influences with system‑wide effects.
    \item Socioeconomic constraints received votes due to their pervasive influence on both access and willingness.
    \item Participants discussed whether risk perception should be treated as a distinct factor or a composite of multiple influences.
    \item Some participants expressed concern about under‑representing structural issues such as transport and clinic placement.
    \item The facilitator recorded which factors received the highest vote totals to inform later modelling prioritisation.
    \item (UNCLEAR) Light side conversation during voting.
    \item Participants broke for lunch.
\end{itemize}

\subsection{Interventions}

\subsubsection{Identifying problematic flows in the system}

\begin{itemize}
    \item Participants reviewed the system diagram and discussed which flows appeared most vulnerable or constrained.
    \item Accessibility to vaccination services was identified as a flow with major structural bottlenecks, including limited clinic hours, inflexible appointment systems, and lack of alternative delivery sites.
    \item Parental willingness to vaccinate was highlighted as a problematic flow due to persistent misinformation, fear of side effects, and general mistrust of institutions.
    \item Participants noted that misinformation spreads rapidly and influences perceptions of risk, which then impacts vaccination uptake.
    \item GP motivation and prioritisation were identified as inconsistent, with some practices actively engaging families and others deprioritising vaccination due to workload or perceived low return on effort.
    \item Participants recognised that cultural norms within certain communities can reduce the flow of individuals choosing vaccination.
    \item Low health literacy was discussed as a barrier that slows down or prevents transitions from “susceptible” to “vaccinated.”
    \item Socioeconomic conditions were identified as constraints affecting multiple flows, especially where transport, employment demands, or childcare responsibilities limit access.
    \item Participants discussed that communication flows between the health system and communities are often insufficient or poorly targeted.
    \item (UNCLEAR) Brief fragmented comment referencing recent outbreaks and responsiveness of different groups.
\end{itemize}

\subsubsection{Identifying potential interventions to improve system behaviour (group task)}

\begin{itemize}
    \item Groups began brainstorming interventions intended to improve vaccination uptake and reduce susceptibility.
    \item Participants proposed increasing availability of community‑based vaccination sites, including pharmacies, health centres, and pop‑up locations.
    \item It was suggested that expanding access points would reduce reliance on GP practices as the sole delivery channel.
    \item Some participants recommended enabling health visitors or other frontline staff to provide vaccinations directly, reducing pressure on GPs.
    \item Participants discussed implementing supplementary vaccination offers, such as home visits or mobile clinics, to reach families with significant access barriers.
    \item Introducing longer‑term, sustained community engagement programmes was proposed to build trust rather than relying on crisis‑driven outreach.
    \item Participants recommended developing campaigns delivered by independent marketing agencies rather than local government, with messaging tailored to community norms.
    \item Using local influencers and trusted messengers, including community leaders, was proposed to counter misinformation and shape positive norms.
    \item Participants suggested creating a large‑scale communication effort that is humorous, relatable, and tailored for social media rather than purely factual or formal.
    \item The idea of making pork‑gelatine‑free vaccine formulations the default across Birmingham was proposed to simplify communication and reduce confusion.
    \item  Some participants noted that supply caps on gelatine‑free vaccines present operational challenges and may require regional or national policy changes.
    \item Participants proposed training staff across health and social care settings in how to hold difficult conversations about vaccination.
    \item It was noted that motivation and beliefs of healthcare workers themselves influence their ability to advocate for vaccination.
    \item Participants suggested conducting fact‑finding work to better understand why some frontline staff decline vaccines, before designing interventions targeting staff beliefs.
    \item Several participants discussed co‑designing interventions with communities to ensure relevance and cultural acceptability.
    \item The idea of creating a network of community “vaccine ambassadors” or champions was proposed to promote vaccination through everyday interactions.
    \item Participants discussed linking vaccination appointments to other services, such as nursery enrolment or early‑years visits, to streamline the process.
    \item Increased use of reminder systems (texts, letters, outreach calls) was proposed to maintain contact with families.
    \item Participants noted that existing communication often does not reach families effectively due to language differences or low literacy.
    \item (UNCLEAR) Short discussion about feasibility of integrating vaccination with school‑based programmes.
\end{itemize}

\subsubsection{Intervention Discussion}

\begin{itemize}
    \item The whole group reconvened to review the proposed interventions.
    \item Community‑based supplementary vaccination offers were widely supported due to strong evidence of convenience improving uptake.
    \item Participants emphasised that pharmacies and health visitors could play especially impactful roles if appropriately resourced.
    \item  Concerns were raised about the feasibility of making gelatine‑free vaccines the default without policy changes or increased supply.
    \item Participants agreed that simplifying the message—removing the need for families to request gelatine‑free options—could significantly reduce confusion.
    \item There was strong support for using independent marketing agencies to deliver culturally relevant, engaging campaigns separate from formal government branding.
    \item Participants noted that humorous or relatable social media content tends to outperform formal health messages in countering misinformation.
    \item It was highlighted that long‑term consistency in engagement is critical; short‑term campaigns during outbreaks are less effective.
    \item The group reaffirmed the importance of vaccine ambassadors and community champions in addressing trust issues within specific communities.
    \item Participants also noted that empowering champions requires careful vetting, support, and collaboration rather than imposing messages.
    \item A discussion took place on training staff in difficult conversations, recognising that many staff feel unprepared to address concerns.
    \item Participants noted that some staff themselves hold doubts or misunderstandings about vaccines, limiting their ability to advocate.
    \item Fact‑finding was encouraged to better understand these sceptical staff segments before designing behaviour‑change interventions.
    \item Concerns were raised about avoiding approaches that could stigmatise staff or communities.
    \item Participants identified structural interventions—such as improved appointment systems and more flexible service hours—as essential for sustained improvement.
    \item Socioeconomic barriers, including transport and time pressure, were highlighted as needing integrated solutions.
    \item Participants concluded that a combined package of interventions, rather than isolated measures, would likely yield the greatest system impact.
    \item (UNCLEAR) Fragment referencing political constraints around mandatory measures.
\end{itemize}

\subsubsection{Ranking proposed interventions using an anonymous online vote}

\begin{itemize}
    \item Participants were asked to scan a QR code linking to an online voting tool.
    \item The facilitator instructed participants to rank all interventions from highest to lowest priority.
    \item The ranking was to reflect the APPEASE criteria, including acceptability, practicability, effectiveness, affordability, spillover effects, and equity.
    \item Participants were encouraged to consider who would “own” each intervention and what system level it would operate at.
    \item Interventions involving gelatine‑free default vaccines, communication with local influencers, and training staff in difficult conversations received substantial discussion.
    \item Participants completed their rankings individually while reviewing displayed options.
    \item The facilitator monitored incoming results and displayed the emerging order.
    \item The top‑ranked interventions included shifting to gelatine‑free default vaccines, improving communication through culturally relevant messaging, and strengthening community leader engagement.
    \item Training health and social care staff in supportive conversations ranked lower but remained important.
    \item Participants noted that even lower‑ranked interventions were still valuable and should not be dismissed.
    \item (UNCLEAR) Brief comments about aligning interventions with upcoming deadlines.
\end{itemize}

\end{document}